\begin{helper}[Parity of the retained defective boundary]
Fix a defective color $\delta$.  Put
\[
 X=V(T),\qquad
 Z=X\setminus\{x^0_{\gamma(\delta)},x^1_{\gamma(\delta)}\},
\]
and let $B$ be the finite set of exact $\delta$-colored boundary-edge
identities of $Z$ which have an incident vertex outside $X$.  Then the
recorded exit $v^0_\delta$ and the first base vertex $y_0$ belong to $Z$,
every vertex of $Z$ is incident with a unique $\delta$-colored edge
identity, $f^0_\delta\in B$, and $|B|$ is odd.
\end{helper}
\begin{proof}
Write $i=\gammaIndex(\delta)$ and $\gamma=\gamma_i$.  By
\noderef{ScheideDefectiveSetup}, $v^0_\delta$ is an exit vertex,
$\gamma$ is missing there, $f^0_\delta$ has color $\delta$ and exact
ends $v^0_\delta,u^0_\delta$, and $u^0_\delta\notin X$.

We first locate the exit.  Expanding \noderef{ScheideExitVertices} and
\noderef{ScheideExitVertex} puts $v^0_\delta$ in $X$.  The exit-covering
field of \noderef{ScheideStructuralConfiguration}, inherited through
\noderef{ScheideReducedConfiguration}, gives an $r$ for which $\gamma_r$
is missing at $v^0_\delta$.  Its block-missing field gives a base vertex
$y_q$ at which the same color is missing.  The base-position equation and
\noderef{TreeSequenceVertexSet} put $y_q$ in $X$.  By
\noderef{ScheideGlobalMaximumOrder}, $X$ is elementary in the sense of
\noderef{ColoredVertexSetElementary}.  If $v^0_\delta\ne y_q$, that
definition makes their two missing-color sets disjoint, contradicting
the fact that the same block color $\gamma_r$ belongs to both.  Hence
$v^0_\delta=y_q$, so $v^0_\delta$ is a base vertex.

The position equations in \noderef{ScheideStructuralConfiguration} put
$y_0$ in a position below $p$, while the selected block vertices occupy
the distinct positions $p+2i$ and $p+2i+1$.  Since $p\ge2$, vertex
injectivity in \noderef{TreeSequence} shows that neither $v^0_\delta$
nor $y_0$ is a selected block vertex.  Both therefore belong to $Z$.

By \noderef{ScheideGlobalMaximumBoundaryFacts}, the defective color
$\delta$ has at least three boundary identities at $X$.  Choose one.
Its color is $\delta$ by \noderef{BoundaryColorEdges}.  Since $X$ is
closed by \noderef{ScheideGlobalMaximumOrder}, $\delta$ cannot be missing
at any vertex of $X$: otherwise the chosen boundary identity contradicts
\noderef{ColoredVertexSetClosed}.  The complementary definitions
\noderef{MissingColors} and \noderef{PresentColors} therefore supply a
$\delta$-colored incident identity at every vertex of $X$, in particular
at every vertex of $Z$.  It is unique.  Indeed, two distinct such
identities would be adjacent by \noderef{AdjacentEdges} and would
contradict \noderef{IsProperEdgeColoring}.  That properness is the
coloring field of the \noderef{KTriple} contained in the
\noderef{GlobalMaximumTriple} field of the structural configuration.

The exact end equation for $f^0_\delta$, the already proved membership
$v^0_\delta\in Z$, and $u^0_\delta\notin X$ show, through
\noderef{Incident}, \noderef{BoundaryEdges}, and
\noderef{BoundaryColorEdges}, that $f^0_\delta\in B$.

It remains to prove parity.  Let $I$ consist of the pairs $(z,f)$ with
$z\in Z$, $f$ incident with $z$, and $f$ colored $\delta$.  The preceding
existence and uniqueness gives a bijection $Z\to I$.  Let $A$ be the set
of $\delta$-colored identities in \noderef{InternalEdges}$(G,Z)$.  Count
$I$ by edge identity.  By \noderef{LooplessMultigraph}, every member of
$A$ has two distinct ends in $Z$ and contributes two incidences.  Every
member of $B$ contributes one incidence.  There are no other nonempty
fibers: the other end of an identity incident with $Z$ lies in $Z$, in
the deleted pair, or outside $X$, and the middle alternative for a
$\delta$-identity is excluded by \noderef{ScheideDefectiveNoCross}.
Consequently
\[
 |Z|=|I|=2|A|+|B|.
\]
The two deleted block vertices are distinct members of $X$, so
$|Z|=|X|-2$.  The order of $X$ is odd by
\noderef{ScheideGlobalMaximumOrder}; therefore $|B|$ is odd, as claimed.
\end{proof}
